\chapter{LITERATURE REVIEW}
\label{ch:literature}
AR and VR are being investigated across an expanding range of spatial tasks, from finding objects in outdoor environments to receiving instructions tied to a surrounding scene. This breadth makes guidance an important interface-design question: an application must help users determine what to attend to, how to act, and when they have reached the intended goal. This review first situates the study among these applications, then examines voice identity and self-similarity, attention guidance, and the gaze-based interaction that supports the present comparison.

\section{AR and VR Applications for Spatial Assistance}
Kumaran et al. compare compass, radar, and in-world arrow aids during outdoor AR search for virtual objects among physical and virtual scene content~\cite{kumaran2023navigation}. Their study illustrates a setting in which assistance must support search while leaving users aware of the surrounding environment. At a different scale, Yang et al.'s AR navigation preprint combines building information with an embodied guide and directional arrows to assist navigation queries~\cite{yang2025embodiednavigation}. These systems extend assistance beyond a front-facing display, where relevant destinations or objects can be distributed around the user.

Other work addresses spatial instructions within a task rather than movement toward a distant destination. Binetti et al. investigate locating out-of-view objects in head-mounted AR, while Zhang et al. study visual and auditory hints during AR search~\cite{binetti2021outofview,zhang2023seehear}. VR offers a complementary setting in which the arrangement of targets, distractors, and initial viewing directions can be controlled and recorded. Chiossi et al.'s comparison of search across physical and virtual content further connects object search with measurements of efficiency and attention~\cite{chiossi2024searching}.

Together, these applications motivate studying how assistance is presented as well as whether it points toward the correct destination. Follow My Voice concentrates on the voice delivering spatial instructions: it compares self-similar and neutral speech within a common VR search application, rather than comparing AR with VR or replacing a navigation interface.

\section{Voice as a Social and Personalized Interface}
Research on machine-generated speech has considered personality, source attribution, and social presence~\cite{nass1993voices,nass2001personality,lee2005socialpresence}. These perspectives motivate treating voice as more than a neutral carrier of words, although they do not imply that every listener will form the same impression or that a social impression necessarily improves task performance.

In this project, personalization takes two forms. In the neutral condition, the participant chooses among existing voices in the supplied catalog; in the self-similar condition, recorded speech is used to create a synthetic voice resembling that participant. Both forms may produce familiarity or liking, but through different routes: preference is expressed directly when the neutral voice is chosen, whereas resemblance is introduced by synthesis and must be confirmed through participants' perceptions. For this reason, the manipulation is better described as a comparison between two forms of personalization than as personalized versus generic voice.

Research on customizing voice-assistant personality has also examined similarity between assistant and user~\cite{snyder2023onevoice}. Personality similarity, however, is conceptually different from acoustic self-resemblance: a listener can prefer a speaking style without hearing it as their own voice, and a close acoustic resemblance can be experienced as uncomfortable. The planned manipulation check therefore asks about perceived voice similarity, while the experience battery separately asks about liking, naturalness, familiarity, and uncanniness.

\section{Self-Similar Voice and Embodied Collaboration}
Self-similar voice provides the empirical motivation for this study, while auditory guidance supplies the setting in which that voice is encountered. Two particularly relevant studies examine resemblance in different roles: an avatar speaking as the player and a virtual collaborator speaking to the participant. The present study extends this question to a guiding voice without a visible agent body.

Kao et al. analyzed 657 participants in an online programming game, comparing self-similar and assigned dissimilar avatar voices while also manipulating pitch~\cite{kao2021selfsimilar}. Self-similar voice was associated with more puzzles completed, longer voluntary play, and higher similarity identification, competence, relatedness, and immersion. Direct differences on the reported intrinsic-motivation subscales were not significant. These findings motivate examining several dimensions of voice experience rather than assuming that resemblance improves enjoyment or every measure of motivation.

The meaning of performance also differs across tasks. In Kao et al.'s game, longer voluntary play can indicate continued engagement, whereas in the present fixed search task a longer trial can indicate difficulty. Follow My Voice will therefore interpret completion time alongside workload, motivation, and reported reliance, without treating speed as a complete account of the value of self-similar voice. The study carries Kao et al.'s research question into a new task rather than predicting that each of their outcomes will replicate in the same direction.

Guo et al. used a 25-participant, within-subject design crossing appearance and voice self-similarity during virtual jigsaw collaboration~\cite{guo2024doppelganger}. Self-similar voices increased likability and believability ratings while also increasing eeriness. Voice did not significantly affect co-presence or perceived intelligence, although appearance affected those measures. These results support keeping favorable and adverse reactions separate and distinguishing voice effects from those of a visible character, although they provide no evidence about search speed.

Both studies compared self-similar voices with assigned alternatives, whereas Follow My Voice compares the clone with a neutral voice selected by the participant~\cite{kao2021selfsimilar,guo2024doppelganger}. The question is therefore whether resemblance to oneself offers a different experience from a voice one has already chosen. Separate measures of perceived resemblance, liking, and willingness to follow instructions allow the study to examine how these responses relate without assuming that they change together.

The present study also requires a distinction between the embodiment of the participant and that of the agent. Participants wear a headset and use a tracked controller, so their interaction involves physical sensing and action, but the guiding agent is not represented by a co-present animated character. Measures developed for an embodied collaborator are therefore reported as adaptations, and their use does not establish that the same factor structure or interpretation holds for a disembodied voice.

Applications of self-voice to daily reminders and emotional self-voice interventions extend this research beyond games and virtual collaboration~\cite{kim2024selfvoice,fang2025emotional}. They motivate asking whether participants can imagine uses for their own synthetic voice, although they do not show that navigation guidance in that voice will be more effective or that participants will experience it as inner speech. The debrief therefore begins with open questions before introducing probes about eeriness, self-talk, and possible applications.

\section{Naturalness, Familiarity, and Uncanniness}
Voice resemblance is not the only acoustic difference that listeners may notice. Synthesis can change articulation, prosody, loudness, duration, or perceived naturalness. Research on vocal uncanniness provides a reason to keep these qualities visible rather than subsume them under similarity~\cite{diel2024uncanny}. A clone that is recognizable but difficult to understand is not equivalent to a recognizable and intelligible clone.

The experiment will use the same semantic library and provider for both conditions. This is an implementation control rather than proof of equal acoustics, because prepared clips can differ in duration and pronunciation even when their text matches. Pilot checks will therefore inspect intelligibility and comfortable playback levels across neutral and cloned voices, and the actual provider, model, voice, and library identifiers will identify the configuration each participant heard.

Custom experience items and the appealing, eerie, and familiar items selected from Guo et al. will remain separate, since combining them would imply a measurement model not established for this study. A participant may report a voice as both familiar and eerie, and preserving the individual item directions keeps that response interpretable.

\section{Workload, Motivation, and Agent Perception}
Because self-resemblance may be both welcome and uncomfortable, the choice of outcome measures is part of the research argument. Task speed describes efficiency, whereas workload, motivation, and agent perception address different aspects of receiving help. The measures reviewed below preserve these distinctions so that the comparison can reveal a tradeoff rather than requiring one voice to rank higher on every dimension.

Zagermann et al. discuss eye tracking as a means of investigating cognitive load in visual computing~\cite{zagermann2016workload}. This work motivates considering gaze-derived measures but does not justify interpreting an unvalidated count of ray hits as workload. The confirmatory plan therefore relies on subjective workload dimensions and task performance, and pupil, fixation, and saccade analyses remain contingent on separate acquisition and event-validation work.

NASA-TLX supplies six workload dimensions~\cite{hart1988nasatlx}; the project's 0--10 response range and unweighted complete mean are documented adaptations. The Intrinsic Motivation Inventory distinguishes interest/enjoyment from related subscales such as effort, pressure, and value~\cite{ryan1982control,mcauley1989imi,csdtIMI}, and reporting the subscales separately preserves that distinction. The instrument appendix specifies the current wording and reverse keys so that scoring can be reproduced.

The agent-perception measures are taken from Guo et al., who attribute them to their original sources~\cite{guo2024doppelganger,biocca2001networked,moussawi2019perception}. Co-presence, perceived intelligence, and anthropomorphism describe participants' attributions, so a high rating does not indicate sentience or correct decisions. Conversely, deterministic validity checks cannot show that a participant found the agent natural. The design therefore requires both kinds of evidence and keeps their meanings distinct.


\section{Visual Search and Attention Guidance in Extended Reality}
To study a guiding voice, the task must give its instructions a meaningful role. Object search provides that role by requiring the participant to connect guidance with a location and then an intended object. The search literature helps specify this setting; it does not supply the answer to the voice comparison.

Search in a surrounding display differs from search on a single front-facing screen because relevant objects may lie outside the visible field. Turning changes what can be inspected, while eye movements refine attention within the visible region. Work on eye and head movements across an extended field of view is therefore relevant to the task geometry~\cite{stein2024eyehead}. In the present study, this motivates recording both head motion and gaze rather than treating one as a complete substitute for the other.

Several approaches to directing attention in immersive environments inform this design. Grogorick et al. address subtle gaze guidance, while HiveFive concerns immersion-preserving attention guidance~\cite{grogorick2017subtle,lange2020hivefive}. Ariza and colleagues investigate vibrotactile assistance and proximity-based multimodal feedback for selection~\cite{ariza2017vibrotactile,ariza2018proximity}. These works identify relevant design alternatives, but the forms of guidance they study are not interchangeable: a cue that changes visual salience, a vibration that indicates proximity, and a spoken spatial instruction impose different perceptual and interpretive demands.

RippleVision provides a recent visual counterpart to gaze-contingent spoken guidance. Kudnick et al. describe brightness ripples restricted to a cone between a point of interest and the current gaze location, with opposite brightness modulation in the two eyes~\cite{kudnick2026ripplevision}. This description is limited to the technique outlined in the published abstract; detailed methods and comparative outcomes are not evaluated here. The technique illustrates how assistance can adapt to a participant's gaze without identifying the target verbally. It also raises the question of how much guidance users find helpful before it becomes intrusive, motivating attention to experience as well as task performance. Follow My Voice holds the assistance policy common across voice conditions to examine self-similar voice; it neither compares guidance modalities nor assumes equivalent effects of visual and spoken cues.

Binetti et al. examine visual and auditory cues for locating out-of-view objects in head-mounted augmented reality~\cite{binetti2021outofview}, and Warden et al. focus on target cue type and location in augmented-reality search~\cite{warden2022visual}. Together, these studies show why target visibility and cue availability belong in the method description. Follow My Voice does not open with a color-and-shape description of the target; instead, the sequence of directions gradually narrows the participant's search.

The hidden-target design also distinguishes this study from known-target search. A participant who already knows the target can continue searching without further assistance, whereas here the participant must obtain enough spatial information to infer which object is intended. The voice may therefore be heard more often, and may matter more, than in a task where it supplements an already known target. This difference calls for a precise task description, although it does not justify expecting stronger voice effects in advance.

\section{Eye Gaze and the Choice of Interaction Modalities}
\subsection{Gaze as Context for Assistance}
Eye tracking provides an observation of where the user is looking that can be compared with the target's position. This makes it useful for adapting guidance as a search unfolds: a broad turn instruction may be relevant while the target is outside the current view, whereas a small vertical correction may be more useful once the user reaches its region. Eye and head movements contribute differently to search across an extended field of view, motivating the collection of both rather than substituting head direction for gaze~\cite{stein2024eyehead}.

Existing gaze-contingent auditory techniques provide precedents for this adaptation. Losing et al. vary a tone with gaze-to-target distance, and Kwok et al. combine initial orientation with finer guidance~\cite{losing2014guiding,kwok2022twostep}. Their desktop studies support investigating feedback tied to ongoing gaze; they do not establish that this headset implementation or its spoken phrases improve performance. EmBARDiment offers a more recent XR example of using gaze context to support an assistant's interaction~\cite{bovo2025embardiment}. The present system applies this rationale to short spatial instructions while keeping the guidance policy common across voices.

\subsection{Speech, Visual Feedback, and Deliberate Selection}
The modalities have complementary roles. Gaze supplies context, speech communicates a direction, visual feedback helps users understand interaction state, and the right-controller trigger provides an explicit selection. Keeping selection separate from looking allows participants to inspect an object without automatically committing to it. Spoken directions also carry the voice-identity cues needed for the self-similar comparison, unlike a nonverbal tone or a visual arrow alone.

This division motivates the apparatus but is not itself an experimental comparison of modalities. Both voice conditions use the same gaze and controller inputs, and neither includes an unguided control. Moreover, a valid gaze intersection does not establish recognition or comprehension: these require behavioral evidence and participant reports. Tracking validity and calibration therefore constrain when gaze can inform assistance, as discussed in the eye-tracking quality section below.

\section{Gaze-Contingent Auditory Guidance}
Auditory guidance is not a single technique. A sound can encode gaze error through loudness, indicate a location through spatial rendering, or communicate an instruction through words. Only the last of these necessarily carries the verbal and vocal identity cues central to this thesis. The following studies inform the shared search design while leaving the self-similar voice comparison open.

Losing et al. tested gaze-contingent sonification in a small desktop study with ten participants: the loudness of a flute tone varied with gaze-to-target distance, reaching silence at the target~\cite{losing2014guiding}. Benefits depended on the search task, with the clearest improvement in the difficult number-search condition. The study supports linking auditory feedback to current gaze but offers no evidence about any particular spoken phrase. The present adaptation replaces the continuously varying tone with short spatial instructions so that participants hear the assigned voice throughout search; the frequency and timing of those instructions are design choices that still require validation rather than properties established by the sonification study.

Kwok et al.'s Two-Step Gaze Guidance evaluated 69 participants using a three-monitor display and eye tracking, rather than a VR headset~\cite{kwok2022twostep}. Initial spatial-audio orientation followed by finer guidance shortened search times relative to the other tested conditions; accuracy did not differ significantly across conditions. Their implementation also buffered transitions between guidance stages. The transferable idea is a staged change in the information a cue supplies: first orient toward a region, then refine within it. Follow My Voice adapts that structure to wall entry and within-wall directions; the cited study does not validate this headset, spoken wording, full-surround geometry, or transition thresholds.

Binetti et al. provide a complementary headset-based example using visual arrows and spatialized alarm sounds to locate out-of-view targets~\cite{binetti2021outofview}. Adding sound mitigated response-time costs when visual cues and targets occupied different depth planes, but did not remove the corresponding accuracy cost. Because the study had no sound-only condition, it offers conditional support for an auditory channel when visual guidance is constrained rather than evidence that speech is superior to visual cues. Its separate response-time and accuracy results also motivate reporting selection errors separately from completion time in the present voice comparison.

Zhang et al.'s 96-participant AR book-search study more directly concerns verbal instructions~\cite{zhang2023seehear}. Its audio hint specified an approximate shelf once at task onset; it was not repeated gaze-contingent speech. In the sub-study without gaze-playback feedback, hint conditions outperformed the control but did not significantly differ from one another. Advantages of combined visual and audio hints over either modality alone appeared in the second task of the feedback sub-study. These results support examining task completion and workload together and stating cue timing and feedback exposure explicitly, although they do not show that repeated spoken guidance or self-similar delivery would be more effective.

Sabouni et al. evaluated nonverbal proximity alerts during a VR search for target paintings~\cite{sabouni2025guidance}. A brief beep signaled entry into a target's vicinity. Compared with no cue, the audio condition reduced task-list checking, but it did not significantly reduce completion time relative to the control. The study therefore supports examining attention allocation alongside completion time, while leaving the effects of spoken directions and self-similar voice unresolved.

Verbal spatial guidance also has accessibility applications. Coughlan et al.'s CamIO system guided a tracked stylus toward points of interest on physical objects using cardinal-direction instructions~\cite{coughlan2020spatial}. Its feasibility study involved four blind participants. This provides a precedent for communicating spatial corrections through speech, but it does not validate gaze-contingent assistance for people with low vision or establish a benefit of self-similar voice. Accessibility extensions of the present system would require separate evaluation with the intended users.

\paragraph{Design synthesis for the voice comparison.}
The resulting design uses a common coarse-to-fine guidance policy to make the two voice conditions comparable. The agent first helps locate the relevant wall, then narrows the search within it; the participant confirms the inferred target with the right trigger. The target's initially hidden identity gives the voice an ongoing instructional role, unlike a single announcement supplementing a known target. Whether hearing these instructions in one's own synthetic voice changes performance, comfort, motivation, or reliance remains the empirical question.

The design records guidance during search, acknowledgments of qualified gaze or controller contact, and feedback after selection as separate events. A success chime after selection cannot explain an earlier change of direction within the same trial, although it may shape later experience. The researcher-facing Trial Review tool should likewise be distinguished from Zhang et al.'s participant-facing gaze-playback intervention~\cite{zhang2023seehear}: reviewing a recording to audit delivery does not expose participants to a learning intervention, and the current protocol includes no participant playback.

The current study does not include an unguided or gaze-unaware control; both voice conditions receive the same gaze-contingent policy. A factorial combination of voice and gaze contingency was considered earlier in the project but is not the design studied here, and its interpretation does not transfer to the two-condition comparison. Any observed difference will therefore concern the assigned voice within the shared assistance system.

\section{Context and Short-Term Memory in XR Agents}
A related AR navigation system by Yang et al. combines Building Information Modeling with retrieval-augmented generation and compares directional arrows alone with an embodied guide accompanied by the same arrows~\cite{yang2025embodiednavigation}. Its comparison concerns the presentation of navigation assistance, not self-similar voice. It provides a useful contrast for the present framing: the question here is how the guiding voice is experienced within a shared assistance mechanism, rather than whether a new reasoning architecture or visible agent improves navigation. The walking task, destination queries, and embodied interface in that system differ from the seated hidden-target task used here.

EmBARDiment uses gaze-informed contextual memory to support interaction with an embodied XR agent~\cite{bovo2025embardiment}. That approach motivates making recent interaction evidence available to the assistance policy, although Follow My Voice implements a narrower, bounded history of recent spatial progress, actual speech, and guidance transitions. Its retention limits are implementation choices rather than thresholds taken from the cited system, and it provides neither long-term autobiographical memory nor an inferred model of the participant's mental state. Unlike the cited interface, the guiding voice also has no visible agent body.

Work on memory-augmented AR agents and the GenLARP system illustrates further directions for context-aware XR agents~\cite{choi2025memory,yu2025genlarp}. The former is a conceptual framework and the latter a short system description, so neither supplies a validated timing threshold for this application. Research on user-aware virtual agents and proactive AR assistance likewise informs the design of state-dependent behavior~\cite{gunawardhana2024useraware,li2025satori}. These works frame the design problem rather than show that this implementation's memory mechanism improves performance.

Short-term history is operationally useful when the same current geometry can have different meanings. A target to the left may be one the participant has not yet approached, or one they have just passed while turning right. A recently delivered cue can make another identical cue unnecessary. History distinguishes these cases only if it records reliable observations and actual playback. Holding this mechanism constant across voices means the study cannot estimate its independent effect.

Feedback specificity offers a further distinction between adaptive assistance and generic encouragement. Zimmerer et al. examined performance-specific feedback from a virtual exercise coach and found changes in selected agent-perception measures without a corresponding effect on motivation~\cite{zimmerer2025feedback}. For Follow My Voice, this distinction motivates describing whether a cue reports observed progress, corrects direction, acknowledges qualified gaze or pointing, or confirms a successful selection. These functions remain common across voice conditions, so the study will compare voice configurations rather than estimate an independent effect of feedback specificity. Results from embodied exercise coaching do not establish an effect on search time in the present task.

\section{Eye-Tracking Quality and Interpretation}
Eye-tracking quality involves more than a nominal sampling frequency. The reporting guidance of Dunn et al. and the data-quality discussion of Niehorster et al. motivate reporting acquisition, validity, spatial performance, and missingness~\cite{dunn2024reporting,niehorster2026dataquality}. The system must establish how many distinct source observations arrive, how old they are when used, and how error varies across the search space.

Validation work on the Vive Pro Eye and tools such as GazeMetrics are methodological resources~\cite{schuetz2022vive,sipatchin2021usability,adhanom2020gazemetrics}. They cannot substitute for validation on the HTC Vive Focus Vision because device model, fit, calibration, coordinate conversion, and task geometry can differ. The protocol therefore requires independent targets across heights, boundaries, distances, and head orientations, and these validation targets must be separate from targets used to estimate an offset correction.

Drewes et al. compare gaze during locomotion in virtual reality and the real world~\cite{drewes2021locomotion}. Because the present study is seated and relies on head and chair rotation rather than free walking, locomotion findings cannot be imported as direct estimates of gaze behavior in this task. Head motion will be logged because it participates in search, while claims about eye events remain bounded by the available measurements.

Salvucci and Goldberg discuss methods for identifying fixations and saccades in eye-tracking data~\cite{salvucci2000fixations}. The 100-ms vicinity qualification used here instead governs interface acknowledgment; it is not a validated fixation classifier. The thesis therefore reports gaze contacts and qualified vicinity events as interaction measures.

\section{Comparison of the Closest Research Strands}
Table~\ref{tab:related-work-map} summarizes what each selected research strand contributes to this design and what it leaves unresolved. The table condenses the preceding review rather than offering a systematic inventory or a ranking of systems. No single strand provides a baseline for the present study, because spatial guidance, voice resemblance, context management, and measurement quality address different questions.

\begin{table}[htbp]\centering
\caption[Related research and boundaries of transfer]{Selected research strands and their role in Follow My Voice. These are methodological connections, not evidence that the present implementation reproduces the cited effects.}
\label{tab:related-work-map}
\small\tagpdfsetup{table/header-rows={1}}
\begin{tabularx}{\linewidth}{@{}p{0.26\linewidth}XX@{}}\toprule
Research strand & Relevance to this project & Boundary of transfer \\\midrule
Self-similar voice~\cite{kao2021selfsimilar,guo2024doppelganger} & Motivates the voice manipulation and both favorable and adverse experience measures. & Educational or embodied collaboration tasks do not establish a search-speed benefit here. \\
Auditory and staged gaze guidance~\cite{losing2014guiding,kwok2022twostep} & Supports investigating gaze-dependent, coarse-to-fine assistance. & Does not validate this headset, phrase policy, or gate durations. \\
AR search and task feedback~\cite{zhang2023seehear} & Separates in-task cues from post-task feedback. & Known-target task and feedback manipulation differ from this hidden-target voice comparison. \\
Context-aware XR agents~\cite{bovo2025embardiment} & Motivates explicit recent context for state-dependent assistance. & Does not establish an independent benefit of this bounded memory. \\
Eye-tracking reporting and validation~\cite{dunn2024reporting,niehorster2026dataquality,adhanom2020gazemetrics} & Informs source timing, validity, spatial validation, and gap reporting. & Reporting guidance and other-device results are not Focus Vision accuracy measurements. \\\bottomrule
\end{tabularx}
\end{table}

Together, these strands frame the particular combination of task and comparison examined here. The agent knows an initially unnamed target, the participant receives adaptive spatial speech, and explicit right-trigger selection remains separate from gaze. The empirical manipulation changes the voice configuration while retaining that assistance mechanism. This combination defines the scope of RQ1--RQ3 without implying that voice familiarity, memory, and gaze contingency have been manipulated independently.

\section{Synthesis and Research Gap}
The selected literature motivates a focused question about self-similar voice: how does guidance spoken in a synthetic version of the participant's own voice compare with guidance spoken in a neutral voice? Prior voice studies make both favorable and uncomfortable reactions plausible, while audio-guidance studies inform a shared task through which those reactions and performance can be examined.

None of the selected work examines this voice comparison during search for an initially unnamed target in a surrounding virtual environment. The present design gives speech an ongoing instructional role and compares the self-similar voice with a neutral voice that already reflects the participant's preference. Its contribution will be to estimate whether and how the two voice configurations differ in this setting, including outcomes on which neither configuration is consistently favored. The next chapter describes the shared apparatus that makes the comparison possible.
